\chapter{Stationarity}

The speeder, settled on sixty, reaches for the cruise control. The needle --- which has held a value, kept a before against an after, made itself carriable, taken on a size, and just now been committed to --- does a new thing: it goes still. The speed holds level. The car is still moving, the engine still working, the road still passing underneath, but the reading sits at one mark and stays there. This chapter is about that stillness --- what it takes for a reading to hold steady, which turns out to be stranger and more exact than a needle simply not moving.

Steady cannot mean that nothing is happening, because plainly everything is. It means something sharper: the conditions can be nudged --- a little more throttle, a slope, a different way of arriving at the same speed --- and the held reading does not move. Some nudges the reading answers to, and some it is simply blind to: change them however you like and the needle does not notice. The chapter's first task is to sort the one from the other --- to find which directions the steady reading reads and which are nothing to it, a difference the reading itself draws by what moves it and what does not.

When a nudge does move the reading, the chapter is not content with roughly how much. It wants the exact change --- the whole of it, signed, with nothing rounded off and nothing left to a limit nobody could finish taking. And it gets it the only honest way: by holding the reading before the nudge against the reading after, and taking the difference. That difference is not an estimate of how the reading shifted; it is the shift itself, read straight off two held values. The chapter makes that exactness its discipline --- the first change a reading shows under a nudge is a thing the apparatus can hold in its hand, not a thing it has to chase.

A reading is steady, then, in a strong and testable sense: nudge it any way the conditions allow, and to first order it does not budge --- every push leaves no trace in the reading. That is balance, and it is the chapter's name for a needle truly held level. But here is the thing to keep hold of, because the chapter turns on it: steady is not flat. A speed held level is not a speed with nothing under it. The needle can be perfectly still and the road beneath it still full of texture --- grades, gusts, the engine's work --- that the first nudge is too coarse to show. The stillness of that first change is not emptiness; it is the apparatus going quiet enough to register what the first order cannot reach.

And whatever the stillness leaves over --- the gap between a reading gone quiet and a reading truly settled --- the chapter finds is not one thing with one name. The same single leftover, read as a balance not yet struck, as a bend that has not quite closed, as a calculation not yet signed off, is one number wearing three names at once. The chapter ends there, on that one residual, held three ways. What it does not yet do is say what lies under the steadiness --- the texture the held needle registers but cannot, at first order, show. That waits for the chapters that come.

\section{The Transparent Middle}

Fix a start $a$ and an end $c$, and route the claim through a middle $b$. Three properties of the middle organize the chapter, and they are genuinely distinct.

The middle is \emph{admissible}, or transparent, when the composed reading equals the direct reading:
\[
  w(a,c) = w(a,b) + w(b,c).
\]
This is the finite triangle relation written as an equality rather than the familiar inequality. A transparent middle bills exactly the through-fare: the detour through $b$ costs what the direct route costs, so the apparatus may tange $b$ into the record without changing what the route claims to cost. The carried claim funges through the transparent middle and arrives with the same receipt it would have carried directly. Whenever the cost never rewards a detour, the inequality $w(a,c) \le w(a,b) + w(b,c)$ holds for every middle; admissibility is the sharp case in which the detour neither saves nor wastes, and that inequality closes to an equality.

The middle is \emph{stationary} when no alternative is cheaper:
\[
  S(b) \le S(b') \qquad \text{for every alternative middle } b',
\]
that is, $w(a,b) + w(b,c) \le w(a,b') + w(b',c)$ for all $b'$. A stationary middle minimizes the composed cost: among every way the apparatus could route the claim through one intermediate mark, $b$ spends the least. Stationarity weighs $b$ only against its alternatives; it says nothing about the direct reading, and a stationary middle may still bill more than $w(a,c)$ when no detour is as cheap as the through-route.

The middle is \emph{flat}, or pure gauge, when the composed cost does not depend on it at all:
\[
  S(b') = S(b) \qquad \text{for every middle } b'.
\]
A flat middle is one the account cannot read. The apparatus carries $b$ in the record, yet the cost it keeps reads the same wherever the middle is placed; no movement of $b$ funges into the count.

These three conditions are not interchangeable, and the section's law states the one implication that always holds: a flat middle is stationary. The argument is immediate. If $S(b') = S(b)$ for every $b'$, then in particular
\[
  S(b) \le S(b') \qquad \text{for every } b',
\]
since a count is never greater than itself. Flatness is the stronger condition: it forces stationarity but says more. A stationary middle sits at the bottom of the composed cost; a flat middle stands on level ground, with no bottom to distinguish it. The converse fails -- a strict minimizer is stationary without being flat -- so the two conditions must be held apart.

The condition the chapter carries forward joins transparency to flatness. Call the middle \emph{closed-flat} when it is both admissible and flat:
\[
  \text{closed-flat}
  \;\equiv\;
  \bigl(\, w(a,c) = w(a,b) + w(b,c) \,\bigr)
  \;\wedge\;
  \bigl(\, S(b') = S(b) \ \text{for all } b' \,\bigr).
\]
A closed-flat middle is transparent and pure gauge at once: it adds nothing to the direct account, and it can be moved anywhere without disturbing the composed account. The residue that closes the chapter is stated for closed-flat middles, because a complete close needs both -- a middle that costs nothing extra and a middle the account cannot feel. Such a middle is the apparatus's first coordinate that it must carry and cannot read: marked in the record, absent from the cost, real enough to name and idle enough to move anywhere. A direction the account cannot feel is a direction its truth cannot be moved by, and the gap between what the apparatus carries and what its cost can register is the whole subject of the chapter.

\section{The Exact First Variation}

The previous section compared one middle with all its alternatives. This section writes that comparison as a signed difference. Keep the start $a$ and the end $c$ fixed, let the route pass first through $b$, and then try another middle $b'$. The apparatus does not change the claim or the endpoints. It changes only the middle mark and asks how the composed reading changes.

The change has two faces. On the left side of the middle, replacing $b$ by $b'$ changes the incoming segment by
\[
  \delta_L(b,b') = w(a,b') - w(a,b) \in \mathbb{Z}.
\]
On the right side, the same replacement changes the outgoing segment by
\[
  \delta_R(b,b') = w(b',c) - w(b,c) \in \mathbb{Z}.
\]
The segment costs themselves live in $\mathbb{N}$, but a comparison of two costs needs signs. A replacement can increase the incoming bill, decrease it, or leave it alone; the same holds on the outgoing side. The signed integers record that direction without pretending that a negative cost has appeared. They measure the change in the bill, not the bill itself.

The first variation of the composed reading adds the two boundary differences:
\[
  \delta S(b,b') = \delta_L(b,b') + \delta_R(b,b').
\]
Expanding the definitions gives the identity that drives the chapter:
\[
\begin{aligned}
  \delta S(b,b')
    &= \bigl(w(a,b') - w(a,b)\bigr)
       + \bigl(w(b',c) - w(b,c)\bigr) \\
    &= \bigl(w(a,b') + w(b',c)\bigr)
       - \bigl(w(a,b) + w(b,c)\bigr) \\
    &= S(b') - S(b).
\end{aligned}
\]
Nothing remains outside this line. The finite first variation is not an approximation to the composed difference; it is the composed difference. If one writes a remainder
\[
  R(b,b') = S(b') - S(b) - \delta S(b,b'),
\]
then the calculation above gives
\[
  R(b,b') = 0
  \qquad \text{for every alternative middle } b'.
\]
This is the finite gift. In the continuum, a derivative usually enters through a limiting Taylor expansion, and the higher-order tail must be estimated, bounded, or made to vanish by a limiting argument. Here the apparatus compares two finite routes directly. The whole change funges into the two boundary faces, so the remainder does not become small. It vanishes exactly.

The same computation can be read as a linearization with no loss. From the base middle $b$, define
\[
  L_b(b') = \delta_L(b,b') + \delta_R(b,b').
\]
Then every allowed replacement of the middle satisfies
\[
  S(b') = S(b) + L_b(b').
\]
The symbol $L_b$ does not introduce a hidden vector space; it names the signed boundary account taken from the chosen base middle. The equation says that the finite apparatus already has the whole first-order account in hand. No small parameter, neighborhood, or limiting process decides whether the approximation is trustworthy. Once the two segment counts exist, the signed boundary account funges the variation into the composed reading with nothing left over.

This exactness also separates two tests that ordinary language easily conflates. The stationary condition from the previous section says that $b$ minimizes the composed reading:
\[
  S(b) \le S(b') \qquad \text{for every } b'.
\]
Using the identity above, the same condition says
\[
  \delta S(b,b') \ge 0 \qquad \text{for every } b'.
\]
No allowed variation lowers the cost. Some variations may raise it; a strict minimum does exactly that for at least one alternative. Stationarity therefore speaks in an inequality.

Flatness speaks in an equality:
\[
  \delta S(b,b') = 0 \qquad \text{for every } b'.
\]
Because $\delta S(b,b') = S(b') - S(b)$, this equality says precisely that every alternative middle carries the same composed cost. Flatness is the pure-gauge case: the account carries the middle but cannot read it. It implies stationarity because zero is nonnegative, but it does more than stationarity. A stationary middle may sit at a bottom; a flat middle erases the bottom by making the whole floor level.

The apparatus now has a first-order detector. It does not need an ambient tangent space, an infinitesimal displacement, or a hidden continuum limit. It takes one held route, varies the middle to another held route, and tanges the signed difference into a receipt. The receipt has two boundary faces, left and right, and their sum is the whole change. When that sum never goes negative, the middle is stationary. When that sum always vanishes, the middle is flat. The next section turns this exact first variation into the balance law at the middle: the two faces must meet as one finite residual. What the continuum reaches only by shrinking a step toward zero, the apparatus holds in full at every step: the exact change in a cost is not a tendency read in the limit but a difference already written on two faces. A measurement that can read how it changes without ever leaving the finite owes its first-order sense to nothing infinite.

\section{Balance at the Middle}

The exact first variation gives the account a local test at the middle. The left boundary contributes one signed change, the right boundary contributes another, and the middle balances only when those two contributions cancel. Write
\[
  r_1(b,b') =
  \delta_L(b,b') + \delta_R(b,b')
  =
  \delta S(b,b').
\]
This number is the first residual of the middle. It measures how much signed cost remains after the incoming and outgoing changes meet. If the left boundary gains what the right boundary loses, or the right boundary gains what the left boundary loses, the two faces close to zero. If their sum does not vanish, the middle has left a first-order trace in the account.

The Euler-Lagrange balance condition is therefore
\[
  r_1(b,b') = 0
  \qquad \text{for every allowed alternative middle } b'.
\]
This is not the same statement as the minimizer condition from the first section. The minimizer condition says
\[
  \delta S(b,b') \ge 0
  \qquad \text{for every } b',
\]
so no replacement lowers the composed cost. The balance condition says
\[
  \delta S(b,b') = 0
  \qquad \text{for every } b',
\]
so every replacement leaves the first variation quiet. A strict minimum may allow positive variations away from $b$; balance forbids even those. The chapter uses \emph{stationary} for the minimizer when the account compares costs, and it uses \emph{balance} or \emph{first-order quiet} for the equality condition at the middle. The distinction keeps the finite calculus honest.

The distinction also fixes the role of the two boundary faces. A minimizer reads the whole composed bill and asks whether any other middle lowers it. It may win that comparison by sitting at the bottom of a shaped cost landscape. Balance asks a sharper question: after the apparatus funges the signed change into left and right contributions, does anything remain at the middle? The answer is local in the finite sense of this chapter. The left and right differences are still computed from complete segment costs, but once they have been computed, their sum is the entire first-order record. No hidden remainder can rescue a nonzero $r_1$, and no limiting argument must be invoked before the apparatus may trust a zero one.

Thus $r_1$ names the middle's first-order residue. It is positive when the replacement raises the composed reading, negative when the replacement lowers it, and zero when the left and right changes cancel exactly. If $r_1(b,b')$ vanishes for every allowed $b'$, the middle has no first-order preference among alternatives. The account can still ask a second question: even if the middle balances against its alternatives, did the route through $b$ cost the same as the direct route from $a$ to $c$? That second question belongs to the direct channel, not to the Euler-Lagrange channel.

The middle also has a direct residual, independent of variation. It compares the direct route with the composed route:
\[
  r_0(b) = w(a,c) - S(b).
\]
Transparency, or admissibility, is exactly the vanishing of this residual:
\[
  r_0(b) = 0.
\]
When $r_0(b)$ vanishes, the detour through $b$ bills the same cost as the direct route. When it does not vanish, the middle may still minimize the composed cost among alternatives, but it has not disappeared from the direct account. The apparatus can then say where the failure lives: the middle may fail at the direct channel, at the first-order channel, or at both.

The two-channel detector combines them:
\[
  R_{\mathrm{tot}}(b,b') = r_0(b) + r_1(b,b').
\]
This total residual does not erase the difference between the channels. It places them in one report. The first term asks whether the detour costs the same as the direct route. The second term asks whether moving the middle leaves the first variation quiet. A total quiet condition requires both:
\[
  r_0(b)=0
  \qquad\text{and}\qquad
  r_1(b,b')=0
  \quad \text{for every } b'.
\]
Under those two equations, the total residual vanishes:
\[
  R_{\mathrm{tot}}(b,b') = 0
  \qquad \text{for every } b'.
\]

Closed-flat middles make exactly this happen. Closed-flat means admissible and flat. Admissibility gives $r_0(b)=0$. Flatness gives $S(b')=S(b)$ for every middle $b'$, and the exact first-variation identity then gives
\[
  r_1(b,b') = \delta S(b,b') = S(b')-S(b)=0
  \qquad \text{for every } b'.
\]
Both detectors go quiet. The direct/composed residual vanishes, and the middle-balance residual vanishes. The apparatus does not merely declare the middle harmless; it records the two equations that make the claim inspectable.

This two-channel form is the reason the balance law can travel through the rest of the chapter. The instrument tanges an abstract equality into two inspectable registers. The first register holds the transparency question, $r_0(b)=0$. The second register holds the balance question, $r_1(b,b')=0$ for every allowed replacement. Once those registers exist, the result becomes fungible: later sections can read it as a direct/composed comparison, as a first-order cancellation, or as a total residual without changing the underlying account. The formula
\[
  R_{\mathrm{tot}}(b,b') = r_0(b) + r_1(b,b')
\]
therefore does not blur the channels. It packages them so the same held record funges through several equivalent readings.

If both answers return zero, the record can pass forward as a residue-free stationarity datum in the broad sense of the chapter: transparent to the direct route and first-order quiet at the middle. If either channel returns a nonzero value, the apparatus has not failed; it has located the residue. A nonzero $r_0$ says the detour still differs from the direct bill. A nonzero $r_1$ says the middle still leaves a first-order trace under some allowed replacement. The next section uses that location to name why the same residual can be read in more than one way without changing the underlying account. Balance is not the absence of a result but a result that has answered twice and stayed quiet: a middle is harmless only when the apparatus can show, in two separate registers, exactly why it costs nothing extra and moves nothing at all. A finite calculus earns its zeros; it never assumes them.

\section{One Residual, Three Names}

The previous section separated the direct channel from the first-order channel. The first-order channel now receives two further readings. Start with the middle-balance residual
\[
  r_1(b,b') = \delta_L(b,b') + \delta_R(b,b').
\]
The same signed number can be read as the residual of a spline closure. Write
\[
  r_{\mathrm{spline}}(b,b') = r_1(b,b').
\]
This equation does not introduce a new failure. It gives the same failure another name. The spline reading asks whether the finite record closes when the middle and its variation are treated as adjacent knots in an account of bending and return. If the residual is nonzero, the closure still carries a visible kink. If the residual is zero for every allowed $b'$, the spline account has no first-order gap at the middle.

The tridiagonal solver reading gives the same residual a computational certificate. Write
\[
  r_{\mathrm{tri}}(b,b')
  =
  r_{\mathrm{spline}}(b,b')
  =
  r_1(b,b').
\]
Here \emph{tridiagonal} names a finite three-neighbor account: previous, current, and next entries carry enough information for the solver certificate to report its residual. The solver does not manufacture balance. It consumes a finite trace whose terminal residual has been solved and identifies that terminal residual with the same residual already present in the gauge and spline readings.

The word \emph{certificate} matters. The apparatus does not watch an infinite process approach a limit. It receives a finite record with three pieces of information: the tridiagonal residual it has solved, the spline residual it claims to match, and the equality that identifies the spline residual with the middle-balance residual. Each piece remains inspectable. The solved trace says zero. The matching equation says the solver residual and the spline residual name the same signed count. The spline-gauge equation says the spline residual and $r_1$ name the same signed count. Chaining those equalities gives the collapse.

A completed certificate therefore supplies
\[
  r_{\mathrm{tri}}(b,b') = 0
  \qquad \text{for every allowed } b'.
\]
Using the residual identifications, the same statement gives
\[
  r_{\mathrm{spline}}(b,b') = 0
  \qquad\text{and}\qquad
  r_1(b,b') = 0
  \qquad \text{for every allowed } b'.
\]
The implication runs through equality of records, not through an analytic convergence theorem. The finite certificate already contains the solved terminal residual. Once the apparatus knows that the tridiagonal residual and the spline residual are the same signed account, zero in one name funges into the others.

This keeps the solver in its proper station. It does not assert that every Lanczos-style procedure converges, or that an arbitrary iteration must eventually find the middle's balance. It asserts a conditional fact about a completed finite witness. If the witness supplies the solved tridiagonal residual and the required residual identifications, then the first-order channel has gone quiet. The force lies in the carried equalities, not in a promise about all possible computations.

This is the point of the three names. In the Euler-Lagrange reading, the residual says whether the left and right boundary changes cancel at the middle. In the spline reading, it says whether the finite closure has a first-order gap. In the tridiagonal reading, it says whether the completed solver trace has a terminal leftover. The names differ because the instrument asks three questions. The count does not differ:
\[
  r_{\mathrm{tri}}(b,b')
  =
  r_{\mathrm{spline}}(b,b')
  =
  r_1(b,b').
\]
One finite residual funges through all three readings without losing its value.

The direct residual remains outside this identification. It still reads
\[
  r_0(b) = w(a,c) - S(b),
\]
and transparency still means
\[
  r_0(b)=0.
\]
The tridiagonal certificate solves the first-order channel. It does not by itself say that the detour through $b$ costs the same as the direct route. When admissibility supplies $r_0(b)=0$ as well, the two-channel report closes:
\[
  R_{\mathrm{tot}}(b,b')
    =
    r_0(b) + r_1(b,b')
    =
    0
  \qquad \text{for every allowed } b'.
\]
Then the residue closes to zero, because the direct channel and the first-order channel have both gone quiet.

This separation prevents a common overread. The first-order certificate can make the middle balanced without making the route transparent. A route may have no first-order gap under variations of the middle and still bill more, or less, than the direct route. The direct/composed account must therefore remain visible. Only when $r_0$ also vanishes does the total residual vanish. The solver certificate contributes the first-order quiet; admissibility contributes the direct quiet.

Closed-flat gives one canonical source of this situation. If the middle is admissible and flat, the direct residual vanishes and the first-order residual vanishes for every allowed replacement. The spline and tridiagonal names then inherit the same zero by identification:
\[
  r_{\mathrm{tri}}(b,b')
  =
  r_{\mathrm{spline}}(b,b')
  =
  r_1(b,b')
  =
  0.
\]
But the section does not need closed-flat as its only route. Its central claim is residual identification: once a finite certificate identifies its tridiagonal terminal residual with the spline residual, and the spline residual with the Euler-Lagrange residual, a solved trace is enough to make the first-order channel quiet.

The truth-terminal case gives the compact grounding example. If every segment cost is zero, then
\[
  w(x,y)=0
  \qquad \text{for all } x,y.
\]
Every route is admissible, every composed reading is flat, and every allowed middle is closed-flat. Therefore
\[
  r_0(b)=0,
  \qquad
  r_1(b,b')=0,
  \qquad
  r_{\mathrm{spline}}(b,b')=0,
  \qquad
  r_{\mathrm{tri}}(b,b')=0.
\]
The example does not replace the residual identification; it displays its least costly instance. At the truth-terminal boundary, the instrument has nothing left to bill, no first-order gap to close, and no solver residue to carry forward.

The three names were never three quantities. The apparatus tanges one signed count and reads it under whichever question the moment asks --- whether two boundary changes cancel, whether a finite closure leaves a gap, whether a solved trace leaves a leftover --- and the count it reads is the same each time. This is the chapter's quiet thesis, and the book's: an invariant is not made many by being named many times. A measurement earns the right to call one residual by three names exactly when no name can change its value, and the whole discipline of a finite calculus is the discipline of keeping that single number answerable under every name it is asked to wear.

\section*{Coda: Stationarity in Review}

Chapter 6 took the committed reading and asked what it takes for it to hold steady. It found the answer in a forced order: a middle the reading might or might not read; an exact first change, held as the difference of two readings rather than chased to a limit; a balance, where every allowed nudge leaves the reading no first-order trace; and a single leftover that turned out to wear three names at once. Drop any one and the next has nothing to test. That is the chapter in review. But set down inside that same order, unnamed, is a sixth part of the thing the chapters have been assembling.

The five parts already on the bench have been, each in its way, unsettled or in motion: a value that can differ from place to place, a value that can differ from before to after, a value read through a floor of disturbance, a value given a size on a finite frame, and a value committed to a definite reading. Now the value can come to rest. It can sit at a reading and hold there --- steady, every allowed nudge leaving it, to first order, unmoved. That is the sixth part: not a value decided once, but a value that stays, balanced against the small pushes around it. Bare, as the others were: not why it rests, not what holds it there, not what its stillness is for --- only that the value can be steady, a reading that does not move to first order. And bare in one more way the chapter was careful about: steady is not empty. Something the first order cannot show still lies beneath the stillness --- but what that is is not this part, and not yet.

What the steadiness is for does not matter right now. Like the five before it, the sixth is set down and left unexplained; what it is for waits on the parts not yet handed over. The book keeps its one habit --- one piece per chapter, in its turn, none before the construction reaches it. Six pieces now lie in order: a value spread across places, the same value moving, the value read through a floor of disturbance, the value sized on a finite frame, the value decided, and now the value at rest. A reader who has watched all six go down may feel the shape pulling tighter, though the book still will not say it.

For now there are six parts, and the sixth is as bare as the rest. It names no rest in particular, fixes no reading, and leans on nothing still to come: a value that can differ from place to place, a value that can differ from before to after, a value read through a floor of disturbance, a value with a size on a finite frame, a value committed to a definite reading, and now a value that can hold steady. It is a sixth mark of a second kind, set down beside the others and waiting, as they wait, for what has not been handed over to give it its place.
